\documentclass[11pt,a4paper]{article}
\usepackage{fontspec,amsmath,amssymb,amsthm,microtype}
\setmainfont{lmroman10-regular.otf}[BoldFont=lmroman10-bold.otf,ItalicFont=lmroman10-italic.otf,BoldItalicFont=lmroman10-bolditalic.otf]
\setmonofont{lmmono10-regular.otf}[Scale=MatchLowercase]
\usepackage[margin=25mm]{geometry}
\usepackage[unicode,hidelinks]{hyperref}
\setlength{\emergencystretch}{3em}
\setlength{\parskip}{0.4em}
\setlength{\parindent}{0pt}
\allowdisplaybreaks[2]
\raggedbottom
\title{\texorpdfstring{Realize every initial coupling independently of the noise}{Realize every initial coupling independently of the noise}}
\author{LangevinLean mathematical exposition}
\date{}
\begin{document}\maketitle
\noindent Source module: \texttt{Unadjusted\allowbreak{}Langevin/\allowbreak{}Initial\allowbreak{}Extension.\allowbreak{}lean}. The accompanying HTML pairs this mathematical exposition with the exact Lean statements and records the scope of the formalization.


Given any coupling \(\gamma\) of \(\mu\) and \(\pi\), extend it to \(\gamma\otimes Q\). This preserves its quadratic cost and supplies independent noise without requiring an optimal coupling.


\section{\texorpdfstring{Coupling.\allowbreak{}is\allowbreak{}Probability\allowbreak{}Measure}{Coupling.isProbabilityMeasure}}


The coupling's first marginal is the initial probability. Pushforward preserves total mass, so \(c(H\times H)=\mu(H)=1\). This constructs normalization of the joint measure using a marginal identity already stored in \texttt{Coupling}.



{\raggedright\noindent Lean declarations: \texttt{Coupling.\allowbreak{}is\allowbreak{}Probability\allowbreak{}Measure}.\par}

\section{\texorpdfstring{initial\allowbreak{}Extension}{initialExtension}}


Extend an initial coupling \(c(dx,dy)\) by an independent noise sample \(s\sim Q\), using \(c\otimes Q\) on \((H\times H)\times S\). Both algorithms retain the prescribed initial joint law while sharing fresh noise independent of the pair.



{\raggedright\noindent Lean declarations: \texttt{initial\allowbreak{}Extension}.\par}

\section{\texorpdfstring{initial\allowbreak{}Extension is\allowbreak{}Probability\allowbreak{}Measure}{initialExtension isProbabilityMeasure}}


The coupling is a probability by its first marginal, and \(Q\) is a probability by hypothesis. Their product therefore has total mass one. This instance supports all subsequent \(L^2\) constructions on the extension.



{\raggedright\noindent Lean declarations: \texttt{initial\allowbreak{}Extension\_\allowbreak{}is\allowbreak{}Probability\allowbreak{}Measure}.\par}

\section{\texorpdfstring{initial\allowbreak{}Extension joint}{initialExtension joint}}


Projecting \(c\otimes Q\) to the first factor integrates out the probability noise coordinate. Since \(Q(S)=1\), the resulting measure is exactly \(c\), preserving the entire initial dependence structure.



{\raggedright\noindent Lean declarations: \texttt{initial\allowbreak{}Extension\_\allowbreak{}joint}.\par}

\section{\texorpdfstring{initial\allowbreak{}Extension initial}{initialExtension initial}}


Compose the joint projection with its first state projection. The extension's joint marginal is \(c\), whose first state marginal is \(\mu\). Pushforward composition therefore yields the specified algorithm initialization law.



{\raggedright\noindent Lean declarations: \texttt{initial\allowbreak{}Extension\_\allowbreak{}initial}.\par}

\section{\texorpdfstring{initial\allowbreak{}Extension reference}{initialExtension reference}}


Compose projection to the initial pair with its second state projection. The coupling's second marginal is \(\pi\), so the reference initialization on the extension has exactly law \(\pi\).



{\raggedright\noindent Lean declarations: \texttt{initial\allowbreak{}Extension\_\allowbreak{}reference}.\par}

\section{\texorpdfstring{initial\allowbreak{}Extension noise}{initialExtension noise}}


The coupling measure has mass one. Projecting its product with \(Q\) to the noise factor therefore returns \(Q\). The \texttt{SFinite} hypothesis supplies the product marginal theorem's measure-theoretic condition.



{\raggedright\noindent Lean declarations: \texttt{initial\allowbreak{}Extension\_\allowbreak{}noise}.\par}

\section{\texorpdfstring{initial\allowbreak{}Extension independent}{initialExtension independent}}


The two coordinate projections of a product probability measure are independent. Here those coordinates are the full coupled initial pair and the fresh noise sample. Independence of either individual initial state from any noise-derived random object follows by composition.



{\raggedright\noindent Lean declarations: \texttt{initial\allowbreak{}Extension\_\allowbreak{}independent}.\par}

\section{\texorpdfstring{initial\allowbreak{}Extension initial noise}{initialExtension initial noise}}


Map \(((x,y),s)\) to \((x,s)\). Product pushforward and the first marginal identity of \(c\) give \(\mu\otimes Q\). Thus removing the reference coordinate leaves the same algorithm input law for every initial coupling.



{\raggedright\noindent Lean declarations: \texttt{initial\allowbreak{}Extension\_\allowbreak{}initial\_\allowbreak{}noise}.\par}

\section{\texorpdfstring{initial\allowbreak{}Extension map}{initialExtension map}}


Factor the map through \(((x,y),s)\mapsto(x,s)\), whose law is \(\mu\otimes Q\). Pushforward composition proves that every measurable output depending only on algorithm initialization and noise has the coupling-independent law \(f_*(\mu\otimes Q)\).



{\raggedright\noindent Lean declarations: \texttt{initial\allowbreak{}Extension\_\allowbreak{}map}.\par}

\section{\texorpdfstring{initial\allowbreak{}Extension chain\allowbreak{}Law}{initialExtension chainLaw}}


Apply \texttt{initial\allowbreak{}Extension\_\allowbreak{}map} to the measurable recursion output. Its canonical pushforward is exactly \texttt{chain\allowbreak{}Law}, so the same output law is realized on every initial-coupling extension. This equality is derived rather than made a field of the coupling.



{\raggedright\noindent Lean declarations: \texttt{initial\allowbreak{}Extension\_\allowbreak{}chain\allowbreak{}Law}.\par}

\section{\texorpdfstring{initial\allowbreak{}Extension iid noise\allowbreak{}Law}{initialExtension iid noiseLaw}}


First project the extension to its product-path noise law, then evaluate coordinate \(n\). The marginal of \(\bigotimes_k\eta\) at that coordinate is \(\eta\), giving the requested noise law on the extension.



{\raggedright\noindent Lean declarations: \texttt{initial\allowbreak{}Extension\_\allowbreak{}iid\_\allowbreak{}noise\allowbreak{}Law}.\par}

\section{\texorpdfstring{initial\allowbreak{}Extension iid independent}{initialExtension iid independent}}


The iid path has independent coordinate evaluations. Pull this independence back along the extension's measure-preserving noise projection. The resulting innovation sequence remains jointly independent on the coupled initialization space.



{\raggedright\noindent Lean declarations: \texttt{initial\allowbreak{}Extension\_\allowbreak{}iid\_\allowbreak{}independent}.\par}

\section{\texorpdfstring{initial\allowbreak{}Extension initial independent}{initialExtension initial independent}}


The algorithm's initial state is a projection of the initial pair. Compose the pair/\allowbreak{}noise independence theorem with this state projection. This proves independence from the entire noise sample, not just from each individual noise vector.



{\raggedright\noindent Lean declarations: \texttt{initial\allowbreak{}Extension\_\allowbreak{}initial\_\allowbreak{}independent}.\par}

\section{\texorpdfstring{initial\allowbreak{}Extension cost}{initialExtension cost}}


The difference \(x-y\) depends only on the initial pair. Its \(L^2\) extended norm under \(c\otimes Q\) equals the same norm under \(c\), because the noise factor has mass one. The latter is exactly \texttt{coupling\allowbreak{}Cost c.\allowbreak{}joint}.



{\raggedright\noindent Lean declarations: \texttt{initial\allowbreak{}Extension\_\allowbreak{}cost}.\par}

\section{\texorpdfstring{initial\allowbreak{}Extension mem\allowbreak{}Lp difference}{initialExtension memLp difference}}


The state-difference function is measurable. Its extended \(L^2\) norm equals the finite initial coupling cost by the previous theorem. These two facts prove \texttt{Mem\allowbreak{}Lp} for the difference on the extension.



{\raggedright\noindent Lean declarations: \texttt{initial\allowbreak{}Extension\_\allowbreak{}mem\allowbreak{}Lp\_\allowbreak{}difference}.\par}

\section{\texorpdfstring{initial\allowbreak{}Extension mem\allowbreak{}Lp reference}{initialExtension memLp reference}}


The reference coordinate has law \(\pi\). Pull its supplied \(L^2(\pi)\) identity moment back through this law to get square integrability on the extension. This establishes the reference initialization without choosing a separate realization.



{\raggedright\noindent Lean declarations: \texttt{initial\allowbreak{}Extension\_\allowbreak{}mem\allowbreak{}Lp\_\allowbreak{}reference}.\par}

\section{\texorpdfstring{initial\allowbreak{}Extension mem\allowbreak{}Lp initial}{initialExtension memLp initial}}


Write \(x=(x-y)+y\). The difference is in \(L^2\) by finite coupling cost and the reference coordinate is in \(L^2\) by its law. Closure of \(L^2\) under addition gives square integrability of the initial algorithm state; no independent moment premise on \(\mu\) is introduced.



{\raggedright\noindent Lean declarations: \texttt{initial\allowbreak{}Extension\_\allowbreak{}mem\allowbreak{}Lp\_\allowbreak{}initial}.\par}

\section{\texorpdfstring{initial\allowbreak{}Extension lp\allowbreak{}Initial}{initialExtension lpInitial}}


Use \texttt{initial\allowbreak{}Extension\_\allowbreak{}mem\allowbreak{}Lp\_\allowbreak{}initial} to convert the first state coordinate into an actual \(L^2\) class. The class is represented almost everywhere by \(((x,y),s)\mapsto x\), preserving its concrete relation to the coupling coordinates.



{\raggedright\noindent Lean declarations: \texttt{initial\allowbreak{}Extension\_\allowbreak{}lp\allowbreak{}Initial}.\par}

\section{\texorpdfstring{initial\allowbreak{}Extension lp\allowbreak{}Reference}{initialExtension lpReference}}


Convert the square-integrable second state coordinate into an \(L^2\) class. Its representative equals \(((x,y),s)\mapsto y\) almost everywhere, with law \(\pi\), and it shares the same product probability space as the algorithm initialization.



{\raggedright\noindent Lean declarations: \texttt{initial\allowbreak{}Extension\_\allowbreak{}lp\allowbreak{}Reference}.\par}

\section{\texorpdfstring{initial\allowbreak{}Extension lp\allowbreak{}Initial law}{initialExtension lpInitial law}}


The representative chosen by \texttt{to\allowbreak{}Lp} agrees almost everywhere with the first state coordinate. Pushforwards are unchanged by such an equality. Combine this with \texttt{initial\allowbreak{}Extension\_\allowbreak{}initial} to obtain the initial law \(\mu\).



{\raggedright\noindent Lean declarations: \texttt{initial\allowbreak{}Extension\_\allowbreak{}lp\allowbreak{}Initial\_\allowbreak{}law}.\par}

\section{\texorpdfstring{initial\allowbreak{}Extension lp\allowbreak{}Reference law}{initialExtension lpReference law}}


Replace the selected \(L^2\) representative by the second state projection almost everywhere. Its already proved marginal is \(\pi\), hence the class used as reference initialization has exactly that law.



{\raggedright\noindent Lean declarations: \texttt{initial\allowbreak{}Extension\_\allowbreak{}lp\allowbreak{}Reference\_\allowbreak{}law}.\par}

\section{\texorpdfstring{initial\allowbreak{}Extension lp cost}{initialExtension lp cost}}


The representative of the class difference equals \(x-y\) almost everywhere. The norm of that \(L^2\) difference is the real value of its extended norm; \texttt{initial\allowbreak{}Extension\_\allowbreak{}cost} identifies it with \((\operatorname{couplingCost}c)_{\mathbb R}\). The finite-cost premise justifies this real conversion.



{\raggedright\noindent Lean declarations: \texttt{initial\allowbreak{}Extension\_\allowbreak{}lp\_\allowbreak{}cost}.\par}

\end{document}
